%% -*- coding: utf-8 -*-
\documentclass[12pt,pagesize,paper=192mm:108mm,landscape]{scrbook} 
%1920x1080 1280x720
\areaset[current]{192mm}{108mm}
\usepackage{calc}
\usepackage[T2A]{fontenc}
\usepackage[utf8]{inputenc}
\usepackage[english,russian]{babel}
\usepackage{microtype}
\usepackage{misccorr}
\usepackage{cmap}
%\usepackage[unicode=true]{hyperref}
\usepackage{graphicx}
\usepackage{amssymb}
\usepackage{amsmath}
%\usepackage{srcltx}
\usepackage{textcomp}
\usepackage{xspace}
%научные символы и смайлики \smiley \frownie
\usepackage{wasysym}
\usepackage{ccicons}
\begin{document}
\begin{titlepage}
  \vspace*{-0.5em}
  \begin{center}    
    % \hspace*{3em}
    % \begin{minipage}[t]{3em}
    %   \includegraphics[width=\textwidth]{../BINP-logo}
    % \end{minipage}\hfill
    % \begin{minipage}{0.23\linewidth}
    % \includegraphics[width=\textwidth]{../NSU-logo}
    % \end{minipage}
    % \hfill
    % \hspace*{6em}

    
    % Кафедра теоретической физики физического факультета НГУ
    % \medskip

    % \Large
    % Профессор Грабовский А.\,В.
    % \smallskip

    % \huge
    % \textbf{Общая теория относительности}
    % \smallskip

    % \Large
    % Лекция № 10
     \vfill

     \normalsize
     \begin{minipage}{0.9\linewidth}
       Пространство время Риндлера как предел пространства времени
       Шварцшильда вблизи горизонта черной дыры. Поверхностная
       гравитация как ускорение статической частицы вблизи горизонта,
       наблюдаемое с пространственной бесконечности. Переход к мнимому
       времени в метрике Риндлера. Условие отсутствия сингулярности у
       евклидовой метрики Риндлера. Периодичность метрики Риндлера по
       мнимому времени. Статистическая сумма канонического
       распределения как след матрицы плотности в энергетическом
       базисе. Вычисление средних значений операторов в
       термостате. Соотношения Кубо"--~Мартина"--~Швингера. Периодичность и
       антипериодичность бозонных и фермионных полей по мнимому
       времени. Выражение статистической суммы через функциональный
       интеграл. Локальная температура вблизи черной дыры "---
       температура Унру. Температура Хокинга как температура черной
       дыры, видимая с пространственной бесконечности. Конформная
       компактификация. Пример: конформная компактификация метрики
       Минковского. Определение пространственной бесконечности $i_0$,
       временной бесконечности будущего $i_{+}$, временной бесконечности
       прошлого $i_{-}$, нулевых бесконечностей будущего $J^{+}$ и прошлого
       $J^{-}$. Статическая вселенная Эйнштейна.
     \end{minipage}
    \vfill

    \normalsize \ccbysa\hspace{0.5em}  Новосибирск 2022
  \end{center}
\end{titlepage}
\end{document}
